\documentclass[11pt,a4paper]{article}

% ── Encoding & fonts ──────────────────────────────────────────────────────
\usepackage[utf8]{inputenc}
\usepackage[T1]{fontenc}
\usepackage[british]{babel}
\usepackage{lmodern}
\usepackage{microtype}

% ── Page layout ───────────────────────────────────────────────────────────
\usepackage[a4paper, top=25mm, bottom=25mm, left=25mm, right=25mm]{geometry}
\usepackage{setspace}
\usepackage{fancyhdr}

% ── Mathematics ───────────────────────────────────────────────────────────
\usepackage{amsmath,amssymb,amsthm,mathtools}

% ── Tables ────────────────────────────────────────────────────────────────
\usepackage{booktabs,longtable,tabularx,array}

% ── Colour & boxes ────────────────────────────────────────────────────────
\usepackage{xcolor}
\usepackage[most]{tcolorbox}

% ── Hyperlinks ────────────────────────────────────────────────────────────
\usepackage{hyperref}

% ── Lists ─────────────────────────────────────────────────────────────────
\usepackage{enumitem}

% ── Hyperref setup ────────────────────────────────────────────────────────
\hypersetup{
  colorlinks = true,
  linkcolor  = blue!60!black,
  urlcolor   = blue!60!black,
  citecolor  = blue!60!black,
  pdftitle   = {UK Regional Housing Market Model},
  pdfauthor  = {Luqman Ahmad},
}

% ── Page style ────────────────────────────────────────────────────────────
\setstretch{1.2}
\setlength{\parskip}{0.6em}
\setlength{\parindent}{0pt}
\setlength{\headheight}{14pt}

\pagestyle{fancy}
\fancyhf{}
\lhead{\small UK Regional Housing Market Model~v0.2.0}
\rhead{\small Luqman Ahmad}
\cfoot{\thepage}

% ── Theorem environments ──────────────────────────────────────────────────
\theoremstyle{plain}
\newtheorem{theorem}{Theorem}[section]
\newtheorem{proposition}[theorem]{Proposition}
\newtheorem{lemma}[theorem]{Lemma}

\theoremstyle{definition}
\newtheorem{definition}[theorem]{Definition}
\newtheorem{remark}[theorem]{Remark}

% ── Custom colour box for governance notes ────────────────────────────────
\tcbuselibrary{breakable}
\newtcolorbox{govbox}[1][Governance Note]{%
  breakable,
  colback   = blue!4!white,
  colframe  = blue!40!black,
  fonttitle = \bfseries\small,
  title     = {#1},
  left=4pt, right=4pt, top=3pt, bottom=3pt,
  sharp corners,
}

% ── Custom macros ─────────────────────────────────────────────────────────
\newcommand{\E}{\mathbb{E}}
\newcommand{\Var}{\operatorname{Var}}
\newcommand{\Cov}{\operatorname{Cov}}
\newcommand{\Pstar}{P^{*}}
\newcommand{\kap}{\kappa}
\newcommand{\sig}{\sigma}
\newcommand{\clip}{\operatorname{clip}}

% ── Title ─────────────────────────────────────────────────────────────────
\title{%
  \textbf{UK Regional Housing Market Model}\\[4pt]
  \large Complete Technical Report --- Version 0.2.0
}
\author{Luqman Ahmad \\ \small Quantitative Analyst}
\date{April 2026}

% ═════════════════════════════════════════════════════════════════════════
\begin{document}

\maketitle
\thispagestyle{empty}
\tableofcontents
\newpage

% ─────────────────────────────────────────────────────────────────────────
\section{Project Overview and Motivation}
% ─────────────────────────────────────────────────────────────────────────

\subsection{Purpose}

This report documents the UK Regional Housing Market Model (v0.2.0), a
research platform estimating equilibrium fair values and stochastic forward
price paths for residential property across 12 UK regions.  The model is
designed for: (a)~systematic research into regional fair value and
mean-reversion dynamics; (b)~scenario-based risk assessment for long-horizon
($\geq$24 month) planning by institutional investors and REIT analysts; and
(c)~comparative regional analysis anchored to Land Registry HPI data.

The model is \emph{not} suitable for individual property valuation,
short-horizon trading decisions, or regulatory capital calculations.  It is
a historical research prototype and does not constitute regulated financial
advice or a stress-testing model as defined under PRA~SS3/18.

\subsection{Seven-Stage Pipeline}

\begin{enumerate}[leftmargin=1.8em, label=\textbf{\arabic*.}]
  \item \textbf{Ingestion} --- fetch and cache raw UK housing and
        macroeconomic data from Land Registry, Bank of England, ONS, and HMRC.
  \item \textbf{Cleaning} --- harmonise, interpolate, and construct the
        canonical monthly panel; record staleness and coverage metadata.
  \item \textbf{Pre-OLS checks and feature engineering} --- multicollinearity
        diagnostics and construction of lagged macro features.
  \item \textbf{OLS calibration} --- semi-structural fair-value layer and
        monthly growth regression; estimation of OU parameters ($\kap$, $\sig$,
        $\gamma$) per region.
  \item \textbf{Monte Carlo simulation} --- 10,000-path Euler-Maruyama
        simulation per region per scenario over a 5-year horizon.
  \item \textbf{Scoring engine} --- consumer and REIT composite scores derived
        from simulation distributional statistics.
  \item \textbf{Streamlit dashboard} --- interactive research UI covering all
        12 English, Welsh, Scottish, and Northern Irish regions.
\end{enumerate}

\subsection{Coverage}

\begin{itemize}[leftmargin=1.8em]
  \item \textbf{Regions:} 12 (all English regions, Wales, Scotland, Northern Ireland).
  \item \textbf{Panel start:} 2005-01-01; \textbf{effective data as of:} 2025-01-01.
  \item \textbf{Simulation horizon:} 5 years (60 monthly steps), 5 macro scenarios.
  \item \textbf{Total paths:} 600,000 (10,000 paths $\times$ 5 scenarios $\times$ 12 regions).
  \item \textbf{Validation status:} APPROVED --- 9/9 sign-off areas complete (2026-04-02).
\end{itemize}

% ─────────────────────────────────────────────────────────────────────────
\section{Mathematical Foundations}
% ─────────────────────────────────────────────────────────────────────────

\subsection{Probability Space and Brownian Motion}

\begin{definition}
A \emph{probability space} is a triple $(\Omega,\mathcal{F},\mathbb{P})$ where
$\Omega$ is the sample space, $\mathcal{F}$ is a $\sigma$-algebra of events, and
$\mathbb{P}:\mathcal{F}\to[0,1]$ is a probability measure with $\mathbb{P}(\Omega)=1$.
\end{definition}

\begin{definition}
A \emph{standard Brownian motion} (Wiener process) $\{W_t\}_{t\geq 0}$ is a
continuous-time stochastic process satisfying:
\begin{enumerate}[leftmargin=1.8em]
  \item $W_0 = 0$ almost surely;
  \item for $0\leq s < t$, the increment $W_t - W_s \sim \mathcal{N}(0,\,t-s)$;
  \item increments over disjoint intervals are independent;
  \item $t\mapsto W_t$ is almost surely continuous.
\end{enumerate}
\end{definition}

Key properties used throughout the model:
\[
\E[W_t] = 0, \qquad \Var(W_t) = t, \qquad \E[W_s\,W_t] = \min(s,t).
\]
The quadratic variation identity $\E[(dW_t)^2] = dt$ is the source of
the second-order correction in It\^{o}'s formula.

\subsection{It\^{o}'s Lemma}

\begin{theorem}[It\^{o}'s Lemma]\label{thm:ito}
Let $X_t$ satisfy the SDE
\begin{equation}
  dX_t = \mu(X_t,t)\,dt + \sig(X_t,t)\,dW_t.
\end{equation}
Let $f:\mathbb{R}\times[0,\infty)\to\mathbb{R}$ be twice continuously
differentiable.  Then
\begin{equation}
  df(X_t,t)
  = \Bigl(
      \frac{\partial f}{\partial t}
    + \mu\,\frac{\partial f}{\partial x}
    + \tfrac{1}{2}\sig^2\,\frac{\partial^2 f}{\partial x^2}
    \Bigr)\,dt
  + \sig\,\frac{\partial f}{\partial x}\,dW_t.
\end{equation}
\end{theorem}

\begin{proof}
Apply a second-order Taylor expansion in $(dx,\,dt)$:
\[
  df = \frac{\partial f}{\partial t}\,dt
     + \frac{\partial f}{\partial x}\,dX_t
     + \frac{1}{2}\,\frac{\partial^2 f}{\partial x^2}\,(dX_t)^2 + \cdots
\]
Substituting $dX_t = \mu\,dt + \sig\,dW_t$ and applying the It\^{o}
multiplication table
\begin{center}
\begin{tabular}{c|cc}
  $\times$ & $dt$ & $dW_t$ \\ \hline
  $dt$     & $0$  & $0$    \\
  $dW_t$   & $0$  & $dt$   \\
\end{tabular}
\end{center}
gives $(dX_t)^2 = \sig^2\,dt + O(dt^{3/2})$.  Collecting terms in $dt$ and $dW_t$:
\[
  df = \Bigl(
         \frac{\partial f}{\partial t}
       + \mu\,\frac{\partial f}{\partial x}
       + \tfrac{1}{2}\sig^2\,\frac{\partial^2 f}{\partial x^2}
       \Bigr)\,dt
     + \sig\,\frac{\partial f}{\partial x}\,dW_t. \qedhere
\]
\end{proof}

\subsection{Geometric Brownian Motion}

Geometric Brownian Motion (GBM) is the classical financial price benchmark:
\begin{equation}
  dS_t = \mu S_t\,dt + \sig S_t\,dW_t.
\end{equation}

\begin{proposition}[GBM Closed-Form Solution]\label{prop:gbm}
The explicit solution of the GBM SDE is
\begin{equation}
  S_t = S_0\exp\!\Bigl[\Bigl(\mu - \tfrac{1}{2}\sig^2\Bigr)t + \sig W_t\Bigr].
\end{equation}
\end{proposition}

\begin{proof}
Apply Theorem~\ref{thm:ito} with $f(S,t) = \log S$.  Then $f_t = 0$,
$f_S = 1/S$, $f_{SS} = -1/S^2$:
\[
  d\log S_t = \Bigl(\mu - \tfrac{1}{2}\sig^2\Bigr)\,dt + \sig\,dW_t.
\]
This is a linear SDE with constant coefficients; integrating from $0$ to $t$
gives $\log(S_t/S_0) = (\mu - \tfrac{1}{2}\sig^2)t + \sig W_t$,
and exponentiating yields the result.
\end{proof}

GBM preserves positivity and is analytically tractable, but it has
\emph{no mean reversion}.  For housing markets, where valuation gaps
attract prices back toward a fundamental anchor, GBM is insufficient
as the primary simulation process.

\subsection{Log Ornstein-Uhlenbeck Process}

\begin{definition}
Let $x_t = \log P_t$ and $m_t = \log \Pstar_t$.  The
\emph{log Ornstein-Uhlenbeck} (log-OU) process is
\begin{equation}\label{eq:logOU}
  dx_t = \kap\,(m_t - x_t)\,dt + \gamma\,dt + \sig\,dW_t,
\end{equation}
where $\kap > 0$ is the mean-reversion speed, $m_t$ is the time-varying
log fair-value target, $\gamma$ is the annual drift, and $\sig > 0$ is
annual volatility.
\end{definition}

\paragraph{Exact solution for constant $m$.}
Rewrite~\eqref{eq:logOU} as
$dx_t + \kap x_t\,dt = (\kap m + \gamma)\,dt + \sig\,dW_t$.
Multiply both sides by the integrating factor $e^{\kap t}$:
\[
  d\bigl(e^{\kap t}\,x_t\bigr) = e^{\kap t}(\kap m + \gamma)\,dt + \sig e^{\kap t}\,dW_t.
\]
Integrating from $0$ to $t$ and dividing by $e^{\kap t}$:
\begin{equation}\label{eq:OUexact}
  x_t = m + \tfrac{\gamma}{\kap}
        + e^{-\kap t}\!\Bigl(x_0 - m - \tfrac{\gamma}{\kap}\Bigr)
        + \sig\int_0^t e^{-\kap(t-s)}\,dW_s.
\end{equation}
The stochastic integral is Gaussian with zero mean and variance
$\sig^2(1-e^{-2\kap t})/(2\kap)$.  The \emph{half-life} of the
log-gap is
\begin{equation}
  t_{1/2} = \frac{\ln 2}{\kap}.
\end{equation}

\begin{remark}
For London ($\kap = 0.200$), $t_{1/2} \approx 3.47$ years.
For Yorkshire and The Humber ($\kap = 0.326$), $t_{1/2} \approx 2.13$ years.
For West Midlands ($\kap = 0.131$), $t_{1/2} \approx 5.29$ years.
Northern Ireland ($\kap = 0.825$) has the fastest reversion:
$t_{1/2} \approx 0.84$ years.
\end{remark}

\subsection{Euler-Maruyama Discretisation}

The model simulates monthly steps using the Euler-Maruyama scheme.
Let $\Delta t = 1/12$ and $x_n = x(n\,\Delta t)$.

\begin{equation}\label{eq:EM}
  x_{n+1}
  = x_n
  + \kap(m_{n+1} - x_n)\,\Delta t
  + \gamma\,\Delta t
  + \sig\sqrt{\Delta t}\;\varepsilon_{n+1},
  \qquad
  \varepsilon_{n+1} \overset{\mathrm{iid}}{\sim} \mathcal{N}(0,1).
\end{equation}

The simulated price is recovered by exponentiation:
$P_{n+1} = \exp(x_{n+1})$.

\begin{remark}[Convergence orders]
The Euler-Maruyama method has strong convergence order $1/2$ and weak
convergence order $1$.  Strong convergence governs pathwise approximation
quality; weak convergence governs the accuracy of expected functional
quantities.  For the distributional statistics (percentiles, probability
of loss) reported in Stage~5, weak convergence at order~1 is the relevant
criterion.
\end{remark}

% ─────────────────────────────────────────────────────────────────────────
\section{Data and Panel Construction}
% ─────────────────────────────────────────────────────────────────────────

\subsection{Data Sources}

\begin{center}
\small
\begin{tabular}{llll}
\toprule
Series & Source & Frequency & Coverage \\
\midrule
Nominal house prices     & HM Land Registry HPI     & Monthly  & 2005--2026 \\
Mortgage rate (2yr fixed) & Bank of England          & Monthly  & 2005--2026 \\
Base rate                & Bank of England          & Monthly  & 2005--2026 \\
Gross earnings (EARN01)  & ONS                      & Annual   & 2005--2025 \\
Private rents            & ONS / Valuation Office   & Monthly  & 2010--2026 \\
CPI                      & ONS                      & Monthly  & 2005--2026 \\
Transactions (SDLT)      & HMRC                     & Monthly  & 2005--2026 \\
Mortgage approvals       & Bank of England          & Monthly  & 2005--2026 \\
Net housing additions    & DLUHC                    & Annual   & 2005--2025 \\
\bottomrule
\end{tabular}
\end{center}

\subsection{Panel Construction}

The canonical panel is built at monthly regional frequency.
Annual series (earnings, net additions) are linearly interpolated to
monthly, and the panel carries explicit staleness flags for all interpolated
columns.  As of the modelling date (2026-04-02), earnings staleness is
13 months across all regions.

Key panel statistics: 2,532 rows $\times$ 41 columns, 12 regions,
starting 2005-01-01.  The fair-value OLS uses \emph{anchor months} only
(June and December of each year), yielding $n = 452$ observations in
the pooled regression.

\subsection{Feature Engineering}

The monthly growth model uses the following constructed features.

\begin{itemize}[leftmargin=1.8em]
  \item \texttt{mortgage\_rate\_lag3}: Bank of England 2-year fixed
        mortgage rate, lagged three months.
  \item \texttt{financial\_stress\_excess\_lag3}: excess of a
        financial-stress composite index above a threshold, lagged three months.
  \item \texttt{earnings\_growth\_12m\_lag1}: 12-month earnings growth,
        lagged one period (annual interpolation; staleness flagged).
  \item \texttt{approvals\_growth}: mortgage approvals 12-month growth.
  \item \texttt{affordability\_gap\_lag1}: price-to-income ratio deviation
        from long-run average, lagged one period.
\end{itemize}

% ─────────────────────────────────────────────────────────────────────────
\section{Fair-Value OLS Model}
% ─────────────────────────────────────────────────────────────────────────

\subsection{Specification}

The fair-value layer estimates equilibrium log real prices using a pooled
OLS regression on anchor months:
\begin{equation}\label{eq:fv}
  \log P^{\mathrm{real}}_{r,t}
  = \alpha
  + \beta_1\log(\text{income}_{r,t})
  + \beta_2\log(\text{rent}_{r,t})
  + \beta_3\;\text{mortgage\_rate}_{r,t}
  + \delta_r
  + \varepsilon_{r,t},
\end{equation}
where $\delta_r$ are regional fixed effects absorbed through regional
dummy variables, and HAC (Newey-West, 6 lags) standard errors correct
inference for autocorrelation.

\textbf{Estimated coefficients} ($n = 452$, $R^2 = 0.9496$):
\begin{align*}
  \hat\beta_1 &= -1.1732 \qquad (\log\text{income}), \\
  \hat\beta_2 &= +1.7443 \qquad (\log\text{rent}), \\
  \hat\beta_3 &= -0.0104 \qquad (\text{mortgage rate, level}).
\end{align*}

The estimated fair value in nominal terms is
\begin{equation}
  \widehat{P}^{\mathrm{nom}}_{r,t}
  = \exp\!\bigl(\widehat{\log P}^{\mathrm{real}}_{r,t}\bigr)
    \times \frac{\mathrm{CPI}_t}{\mathrm{CPI}_{\mathrm{base}}}.
\end{equation}

The log valuation gap is $x_{r,t} = \log P^{\mathrm{real}}_{r,t}
- \log\widehat{P}^{\mathrm{real}}_{r,t}$.

\subsection{Gauss-Markov Theorem}

\begin{theorem}[Gauss-Markov]\label{thm:gm}
Under the classical linear model assumptions---exogeneity
($\E[\varepsilon \mid X] = 0$), homoskedasticity
($\Var(\varepsilon \mid X) = \sig^2 I$), and no perfect
multicollinearity---the OLS estimator
$\hat\beta = (X^\top X)^{-1}X^\top y$
is the Best Linear Unbiased Estimator (BLUE).
\end{theorem}

\begin{proof}
\textbf{Unbiasedness.}
\[
  \E[\hat\beta]
  = (X^\top X)^{-1}X^\top\E[y]
  = (X^\top X)^{-1}X^\top X\beta
  = \beta.
\]

\textbf{Minimum variance (BLUE).}
Let $\tilde\beta = Cy$ be any other linear unbiased estimator.
Write $C = (X^\top X)^{-1}X^\top + D$ for some matrix $D$.
Unbiasedness requires $\E[\tilde\beta] = \beta$, hence $DX = 0$.
The variance of $\tilde\beta$ is
\[
  \Var(\tilde\beta)
  = \sig^2 CC^\top
  = \sig^2\!\Bigl[(X^\top X)^{-1} + DD^\top\Bigr].
\]
Since $DD^\top$ is positive semi-definite,
$\Var(\tilde\beta) \succeq \sig^2(X^\top X)^{-1} = \Var(\hat\beta)$
for every linear unbiased estimator $\tilde\beta$.
\end{proof}

\begin{govbox}[HAC Correction and Gauss-Markov Violation]
Housing price residuals violate the Gauss-Markov homoskedasticity
assumption: all 12 regions exhibit significant autocorrelation and
ARCH-type conditional heteroskedasticity (Section~\ref{sec:validation}).
HAC (Newey-West, 6 lags) standard errors restore asymptotically valid
inference by adjusting $\Var(\hat\beta)$ for the observed autocorrelation
and heteroskedasticity structure.  Point estimates $\hat\beta$ remain
consistent; only standard errors change.  Residual autocorrelation in
housing is well-documented \cite{case1989} and is the primary motivation
for HAC throughout.
\end{govbox}

\subsection{PTI/PTR Blend Weight}

The fair-value blend weight $w_{\mathrm{PTI}}$ allocates between a
price-to-income (PTI) channel and a price-to-rent (PTR) channel.
The weight is estimated by grid-search over
$w \in [0.25,\,0.75]$ minimising pooled squared log-price deviation:
\begin{equation}
  \hat{w}_{\mathrm{PTI}}
  = \arg\min_{w\in[0.25,\,0.75]}
    \sum_{r,t}
    \bigl(\log P_{r,t}
          - w\,\log\text{income}_{r,t}
          - (1-w)\,\log\text{rent}_{r,t}\bigr)^2.
\end{equation}
The estimated value is $\hat{w}_{\mathrm{PTI}} = 0.25$, clipped to the
lower bound, indicating that the rental parity channel dominates in this
panel---consistent with the large estimated $\hat\beta_2 = 1.744$
on $\log\text{rent}$.

\subsection{Mean-Reversion Speed Estimation}

The mean-reversion speed $\kap_r$ is estimated from an AR(1) regression
on the log valuation gap:
\begin{equation}
  x_{r,t} = \phi_r\,x_{r,t-1} + \eta_{r,t}.
\end{equation}
The annualised speed is $\kap_r = -12\log\hat\phi_r$.  Regional
estimates range from $\kap = 0.131$ (West Midlands) to
$\kap = 0.825$ (Northern Ireland).

\subsection{Volatility Estimation}

Annual volatility $\sig_r$ is blended from raw and winsorised
realised monthly volatility:
\begin{equation}
  \sig^{\mathrm{blend}}_r
  = \clip\!\bigl(
      0.35\,\sig^{\mathrm{raw}}_r
    + 0.65\,\sig^{\mathrm{win}}_r,\;
      0.025,\; 0.12
    \bigr).
\end{equation}
An IQR prior raises $\sig_r$ when the implied 5-year distribution
is too narrow relative to historical 5-year return dispersion
(threshold: IQR ratio $\geq 0.30$).
A sigma floor override is applied to South West ($\sig \to 0.045$,
raised from the fitted 0.030) to enforce the IQR ratio requirement.

\subsection{Calibrated SDE Parameters by Region}

\begin{table}[htbp]
\caption{Region-level SDE parameters, bank-grade canonical rebuild
(2026-04-02).  $\kap$ = annualised mean-reversion speed;
$\sig$ = annual log-price volatility;
$t_{1/2} = \ln 2 / \kap$ in years.}
\label{tab:sde}
\centering\small
\begin{tabular}{lrrrl}
\toprule
Region & $\kap$ & $\sig$ & $t_{1/2}$ (yr) & Note \\
\midrule
East Midlands             & 0.419 & 0.080 & 1.65 & \\
East of England           & 0.218 & 0.100 & 3.18 & \\
London                    & 0.200 & 0.051 & 3.47 & rate $\approx 0$ in OLS \\
North East                & 0.497 & 0.120 & 1.40 & highest $\sig$ \\
North West                & 0.314 & 0.095 & 2.21 & \\
Northern Ireland          & 0.825 & 0.115 & 0.84 & fastest mean reversion \\
Scotland                  & 0.468 & 0.044 & 1.48 & \\
South East                & 0.227 & 0.085 & 3.05 & \\
South West                & 0.458 & 0.030 & 1.51 & floor $\to 0.045$ applied \\
Wales                     & 0.214 & 0.040 & 3.24 & \\
West Midlands             & 0.131 & 0.055 & 5.29 & slowest mean reversion \\
Yorkshire and The Humber  & 0.326 & 0.050 & 2.13 & \\
\bottomrule
\end{tabular}
\end{table}

% ─────────────────────────────────────────────────────────────────────────
\section{Stochastic Simulation}
% ─────────────────────────────────────────────────────────────────────────

\subsection{GBM vs Log-OU: Conceptual Comparison}

\begin{center}
\begin{tabular}{lll}
\toprule
Property & GBM & Log-OU \\
\midrule
Mean reversion     & No         & Yes ($\kap > 0$) \\
Long-run equilibrium & None     & Fair-value path $m_t$ \\
Log-price stationary distribution & None & $\mathcal{N}(m, \sig^2/(2\kap))$ \\
Positivity         & Yes        & Yes (log-price space) \\
Role in this model & Benchmark  & Primary process \\
\bottomrule
\end{tabular}
\end{center}

\subsection{Target Mean Path Construction}

The simulation mean path $m_t$ is constructed as a linear blend
between the current fair-value anchor $\Pstar_0$ and the scenario
terminal target $\Pstar_T$:
\begin{equation}
  m_t = (1 - s_t)\log\Pstar_0 + s_t\log\Pstar_T,
  \qquad s_t = \frac{t}{T}.
\end{equation}
The scenario terminal target is formed from the baseline five-year
log return adjusted by macro overlays:
\begin{equation}
  \log R^{(5\mathrm{y})}_{\mathrm{scen}}
  = \log R^{(5\mathrm{y})}_{\mathrm{base}}
    + 5 \times
      \bigl(
        3\hat\beta_{\mathrm{rate}}\,\Delta r
        + 12\hat\beta_{\mathrm{stress}}\,\Delta f
        + \Delta\gamma
      \bigr),
\end{equation}
where $\Delta r$ is the scenario rate shift (pp), $\Delta f$ is the
stress overlay, and $\Delta\gamma$ is the additional drift shift.

\subsection{Five Macro Scenarios}

\begin{center}
\begin{tabular}{lcc}
\toprule
Scenario & Rate shift (pp) & Drift shift (pp yr$^{-1}$) \\
\midrule
Baseline      & $\phantom{+}0.0$ & $\phantom{+}0.0$ \\
Soft Landing  & $-0.5$ & $+0.5$ \\
Rate Shock    & $+2.0$ & $-1.5$ \\
Stagflation   & $+1.5$ & $-2.0$ \\
Recovery Boom & $-1.0$ & $+2.0$ \\
\bottomrule
\end{tabular}
\end{center}

\subsection{Simulation Procedure}

For each of the 12 regions and 5 scenarios, the Euler-Maruyama
scheme~\eqref{eq:EM} is applied with:
\begin{itemize}[leftmargin=1.8em]
  \item $N = 10{,}000$ independent paths;
  \item $T = 60$ monthly steps ($\Delta t = 1/12$);
  \item i.i.d.\ standard normal draws $\varepsilon_{n+1}$ from a fixed
        random seed for reproducibility.
\end{itemize}

Output statistics per region-scenario: median 5-year growth,
$p_{10}$/$p_{25}$/$p_{75}$/$p_{90}$ percentiles, and the probability
of a terminal loss exceeding 10\%.

% ─────────────────────────────────────────────────────────────────────────
\section{Scoring Engine}
% ─────────────────────────────────────────────────────────────────────────

\subsection{Signal Components}

Let $g_r$ denote the current percentage valuation gap
(positive = overvalued relative to $\Pstar$).  Define:
\begin{align}
  C1_{\mathrm{mispricing}}
  &= \clip\!\bigl(50 - 0.95\times g_r,\;0,\;100\bigr), \\
  C2_{\mathrm{consumer}}
  &= \clip\!\bigl(50 + 1.8\times\text{WeightedConsumerReturn},\;0,\;100\bigr), \\
  C2_{\mathrm{REIT}}
  &= \clip\!\bigl(50 + 1.8\times\text{WeightedREITReturn},\;0,\;100\bigr), \\
  \text{YieldScore}
  &= \clip\!\bigl(50 + 16(\text{GrossYield} - 5.5),\;0,\;100\bigr), \\
  \text{TailScore}
  &= \clip\!\bigl(100(1 - p_{\mathrm{loss}>10\%}),\;0,\;100\bigr).
\end{align}

\subsection{Composite Scores}

\begin{align}
\text{ConsumerScore}
&= \clip\!\bigl(
    0.40\,C1_{\mathrm{mispricing}}
  + 0.35\,C2_{\mathrm{consumer}}
  + 0.25\,\text{TailScore},\;
  0,\;100
  \bigr), \\
\text{REITScore}
&= \clip\!\bigl(
    0.20\,C1_{\mathrm{mispricing}}
  + 0.35\,C2_{\mathrm{REIT}}
  + 0.25\,\text{YieldScore}
  + 0.20\,\text{TailScore},\;
  0,\;100
  \bigr).
\end{align}

\subsection{Scenario Aggregation Weights}

Consumer aggregation is downside-skewed; REIT aggregation is
equal-weighted:
\begin{align*}
  w^{\mathrm{consumer}}
  &= (0.40,\;0.25,\;0.20,\;0.10,\;0.05)
  \quad\text{(Baseline, Rate Shock, Stagflation, Soft Landing, Boom)}, \\
  w^{\mathrm{REIT}}
  &= (0.20,\;0.20,\;0.20,\;0.20,\;0.20).
\end{align*}

\subsection{Score Labels and Governance}

\begin{center}
\begin{tabular}{lll}
\toprule
Score range & Consumer label & REIT label \\
\midrule
$\geq 62$ & Supportive & Supportive \\
$50$--$61$ & Mixed      & Mixed      \\
$< 50$    & Cautious   & Cautious   \\
\bottomrule
\end{tabular}
\end{center}

\begin{govbox}[Score Label Change --- Risk Register R004 (RESOLVED)]
Labels were changed from ``Buy / Hold / Don't Buy'' to ``Supportive /
Mixed / Cautious'' to reflect their nature as descriptive research
buckets, not investment recommendations.  The dashboard explicitly states
that labels do not constitute investment advice.  This change is documented
in \texttt{SCORING\_BACKTEST\_NOTE.md} and risk register item~R004,
resolved 2026-04-02.
\end{govbox}

% ─────────────────────────────────────────────────────────────────────────
\section{Validation and Backtesting}\label{sec:validation}
% ─────────────────────────────────────────────────────────────────────────

\subsection{OLS Replay Deviation}

The canonical rebuild re-estimates the fair-value OLS from the panel
and compares fitted values against deployed outputs:
\[
  \max_{r,t}\,\bigl|
    \log\widehat{P}^{\mathrm{deployed}}_{r,t}
    - \log\widehat{P}^{\mathrm{replay}}_{r,t}
  \bigr| = 0.000
\]
(to machine precision), confirming that all deployed fair-value outputs
are exactly reproducible from source data and code.

\subsection{Subsample OLS Validation}

Training set: observations before 2018-01-01.
Holdout: 2018-01-01 to 2026-01-01 ($n = 1{,}079$).

\begin{center}
\begin{tabular}{lcc}
\toprule
Model & RMSE (1m growth) & Dir.\ accuracy (12m) \\
\midrule
OLS growth model      & 0.02165 & 64.7\% \\
Zero-growth benchmark & 0.01526 & --- \\
\bottomrule
\end{tabular}
\end{center}

\begin{govbox}[Short-Horizon RMSE Disadvantage]
The OLS growth model does not beat the zero-growth benchmark at 1-month
RMSE.  This is expected and consistent with weak-form efficiency in
liquid residential markets: monthly price changes are largely
uninformative.  The model's value-add is at horizons
$\geq 24$ months, where mean reversion dominates.  Directional accuracy
of 64.7\% at the 12-month horizon exceeds the 55\% governance threshold
and is the primary evidence of model skill.
\end{govbox}

\subsection{Sensitivity Analysis (One-At-a-Time Shocks)}

Shocks applied to the fair-value OLS specification~\eqref{eq:fv}.
Responses are uniform across all 12 regions because the OLS
specification uses pooled coefficients.

\begin{center}
\begin{tabular}{lcccc}
\toprule
Shock & Expected sign & Average $\Delta P^*$ (\%) & Sign correct? \\
\midrule
Income $+10\%$  & $+$ & $-10.58$ & NO \\
Income $-10\%$  & $-$ & $+13.16$ & NO \\
Rates $+100$~bp & $-$ & $-0.01$  & YES \\
Rates $-100$~bp & $+$ & $+0.01$  & YES \\
Rent $+10\%$    & $+$ & $+18.09$ & YES \\
Rent $-10\%$    & $-$ & $-16.79$ & YES \\
\bottomrule
\end{tabular}
\end{center}

\begin{govbox}[Income Sign Violation (Frisch-Waugh-Lovell Artefact)]
The OLS income coefficient $\hat\beta_1 = -1.173$ is negative,
producing sign violations on income shocks.
By the Frisch-Waugh-Lovell theorem, the partial regression coefficient
on $\log\text{income}$ in a specification that \emph{already includes}
$\log\text{rent}$ absorbs only the residual variation in income after
partialling out rent.  Since income and rent are positively correlated
cross-sectionally, this residual carries negative partial correlation
with prices in the panel.  This is an econometric artefact of the
semi-structural design, not a model defect.  The income variable is
retained for specification consistency with the broader literature.
Rent ($+18\%$ per $+10\%$ rent shock) and rate channels produce
economically correct signs and magnitudes.
\end{govbox}

\subsection{Confidence Interval Coverage}

A rolling 1-step-ahead test checks whether actual log prices fall
within the model's 90\% predictive interval:
\[
  \bigl|x_{t+1} - \E[x_{t+1}\mid x_t]\bigr|
  \leq 1.645\,\sig\,\sqrt{\Delta t},
\]
where $\E[x_{t+1}\mid x_t] = x_t e^{-\kap\Delta t} + \gamma\,\Delta t$
and the time-varying fair value $m_{t+1}$ is taken from the in-sample
OLS fair-value panel.  Holdout: 2018-01-01 to 2026-01-01
(96 months per region).

\begin{longtable}{lcc}
\caption{Rolling 1-step-ahead 90\% CI coverage by region
(holdout 2018--2026).
Target $= 0.90$;
pass threshold $\geq 0.80$;
material threshold $= 0.70$.}\label{tab:ci}\\
\toprule
Region & Coverage & Assessment \\
\midrule
\endfirsthead
\multicolumn{3}{l}{\small\itshape(Continued from previous page)}\\
\toprule
Region & Coverage & Assessment \\
\midrule
\endhead
\midrule
\multicolumn{3}{r}{\small\itshape(Continued on next page)}\\
\endfoot
\bottomrule
\endlastfoot
East Midlands             & 0.854 & Acceptable \\
East of England           & 0.865 & Acceptable \\
London                    & 0.740 & Monitor \\
North East                & 0.854 & Acceptable \\
North West                & 0.854 & Acceptable \\
Northern Ireland          & 0.885 & Acceptable \\
Scotland                  & 0.698 & Monitor \\
South East                & 0.875 & Acceptable \\
South West                & 0.594 & Below material threshold \\
Wales                     & 0.677 & Monitor \\
West Midlands             & 0.823 & Acceptable \\
Yorkshire and The Humber  & 0.729 & Monitor \\
\midrule
\textbf{Pooled (mean)}    & \textbf{0.787} & \textbf{MARGINAL} \\
\end{longtable}

\begin{govbox}[CI Coverage Finding]
Pooled coverage of 0.787 lies between the 0.70 material threshold and
the 0.80 pass threshold.  South West (0.594) and Scotland (0.698) are
the weakest regions.  Recalibrating $\sig$ to inflate coverage would
constitute grade-shopping and is not recommended; the sigma floor override
for South West (raised to 0.045) addresses the root cause at the parameter
level.  Monitoring is recommended across all four flagged regions.
\end{govbox}

\subsection{Residual Diagnostics}\label{sec:resid}

Fair-value OLS residuals are tested per region using:
Ljung-Box (lags 1, 6, 12) for autocorrelation;
Engle's ARCH LM test (lag 4) for conditional heteroskedasticity;
and the Jarque-Bera test for non-normality.
Statistics in bold have $p < 0.05$.

\begin{longtable}{lrrrrrr}
\caption{Residual diagnostics (fair-value OLS, anchor months).
LB$_k$ = Ljung-Box statistic at lag $k$;
ARCH = Engle ARCH-LM statistic (lag 4);
JB = Jarque-Bera statistic.}\label{tab:diag}\\
\toprule
Region & $n$ & LB$_1$ & LB$_6$ & LB$_{12}$ & ARCH & JB \\
\midrule
\endfirsthead
\multicolumn{7}{l}{\small\itshape(Continued from previous page)}\\
\toprule
Region & $n$ & LB$_1$ & LB$_6$ & LB$_{12}$ & ARCH & JB \\
\midrule
\endhead
\midrule
\multicolumn{7}{r}{\small\itshape(Continued on next page)}\\
\endfoot
\bottomrule
\endlastfoot
East Midlands            & 41 & \textbf{31.9} & \textbf{119.9} & \textbf{142.7} & \textbf{34.3} & 0.27 \\
East of England          & 41 & \textbf{35.6} & \textbf{143.4} & \textbf{182.6} & \textbf{33.9} & 2.48 \\
London                   & 41 & \textbf{33.6} & \textbf{134.9} & \textbf{170.1} & \textbf{34.4} & 0.84 \\
North East               & 41 & \textbf{36.7} & \textbf{139.4} & \textbf{164.7} & \textbf{35.7} & 1.45 \\
North West               & 41 & \textbf{35.3} & \textbf{142.7} & \textbf{171.9} & \textbf{32.5} & 1.31 \\
Northern Ireland         & 21 & \textbf{14.8} & \textbf{35.0}  & \textbf{42.0}  & \textbf{16.4} & 1.85 \\
Scotland                 & 29 & \textbf{22.1} & \textbf{54.1}  & \textbf{55.6}  & \textbf{22.9} & 2.86 \\
South East               & 41 & \textbf{33.7} & \textbf{130.0} & \textbf{164.4} & \textbf{35.0} & 1.40 \\
South West               & 41 & \textbf{33.2} & \textbf{121.6} & \textbf{142.2} & \textbf{35.8} & 0.44 \\
Wales                    & 33 & \textbf{28.3} & \textbf{103.0} & \textbf{113.9} & \textbf{28.0} & 2.19 \\
West Midlands            & 41 & \textbf{35.1} & \textbf{144.2} & \textbf{178.3} & \textbf{33.1} & 1.41 \\
Yorkshire and The Humber & 41 & \textbf{35.0} & \textbf{140.0} & \textbf{169.8} & \textbf{35.2} & 1.13 \\
\end{longtable}

All 12 regions flag autocorrelation at all three lags and ARCH-type
conditional heteroskedasticity.  Zero regions fail the Jarque-Bera
normality test.  Residual autocorrelation in housing price series is
well-documented \cite{case1989} and is the primary motivation for
HAC standard errors throughout.  These diagnostic flags confirm that
HAC correction is necessary and in place; they are not model defects.

\subsection{Model Horse Race}

Four models are compared on holdout 2018-01-01 to 2026-01-01
(pooled across 12 regions, 96 observations per region).
GBM and random walk use 500 simulated paths; the median path is
the point forecast.  Log-OU uses the analytical conditional expectation.

\begin{center}
\begin{tabular}{lcccc}
\toprule
Model        & RMSE (1m)  & RMSE (12m)         & Dir.\ Acc.\ (1m) & Dir.\ Acc.\ (12m) \\
\midrule
Log-OU       & 0.01445    & \textbf{0.04442}   & 0.554            & \textbf{0.806} \\
GBM          & 0.01449    & 0.04881            & 0.575            & 0.790 \\
Zero-growth  & 0.01466    & 0.05582            & 0.055            & 0.000 \\
Random walk  & 0.01472    & 0.05612            & 0.448            & 0.501 \\
\bottomrule
\end{tabular}
\end{center}

\textbf{Key findings:}
\begin{enumerate}[leftmargin=1.8em]
  \item At the 12-month horizon, Log-OU achieves the lowest RMSE
        (0.04442 vs GBM 0.04881 and zero-growth 0.05582) and the
        highest directional accuracy (80.6\% vs 79.0\% for GBM).
  \item At the 1-month horizon, zero-growth has the lowest RMSE,
        consistent with weak-form efficiency---expected and documented.
  \item Log-OU outperforms all three benchmarks at the 12-month horizon,
        demonstrating that the OU mean-reversion structure adds measurable
        value at planning-relevant horizons.
\end{enumerate}

% ─────────────────────────────────────────────────────────────────────────
\section{Dashboard: UI Section Guide}
% ─────────────────────────────────────────────────────────────────────────

The Streamlit research dashboard is structured as eight navigable pages.

\subsection{Page 1 --- Research Overview}

Displays the model validation status, effective data date, and a summary
scorecard across all 12 regions.  Consumer and REIT scores with
Supportive/Mixed/Cautious labels are shown in a regional summary table.
An explicit disclosure banner states that model output is for research
purposes only and does not constitute investment advice.

\subsection{Page 2 --- Regional Signals}

Detailed view for a selected region.  Shows: the current fair-value gap
($x_r$) and a price vs $\Pstar$ chart; score decomposition with
labelled C1/C2/Yield/Tail components; and evidence-level classification
(Validated/Partial/Contextual) for each signal component.

\subsection{Page 3 --- Region Comparison}

Side-by-side bar charts of consumer and REIT scores across all 12 regions.
Supports ranking by score, valuation gap, or directional accuracy.
Designed to support relative regional positioning for portfolio allocation.

\subsection{Page 4 --- Scenario and Diagnostics}

Shows the five macro scenarios for a selected region: fan charts
($p_{10}$/$p_{25}$/median/$p_{75}$/$p_{90}$) of simulated price paths
and terminal return distributions.  Contains the interactive Scenario Lab,
labelled explicitly as a ``coefficient-based perturbation tool'' and not
a full econometric re-run.

\subsection{Page 5 --- Household Explorer}

Interactive tool allowing a user to input a purchase price, deposit, and
household income to see personalised affordability metrics and how current
model signals apply to their budget range.  Carries an explicit disclaimer
that outputs are for illustrative research purposes only.

\subsection{Page 6 --- Methodology}

Plain-language documentation of the seven-stage pipeline, fair-value OLS
specification, simulation setup, and scoring weights.  Points to this
technical report for full mathematical derivations and governance details.

\subsection{Page 7 --- Technical Controls}

Exposes the effective data date, earnings staleness months, model
validation status, simulation seed, and per-region SDE parameters.
Allows toggling between Baseline and alternative scenario views.

\subsection{Page 8 --- Limitations}

Dedicated page listing all material limitations from
\texttt{MODEL\_LIMITATIONS.md}: regional-average vs.\ individual-property
scope; interpolated annual series; omitted-variable risk; Scenario Lab
limitations; and score compression.  This page is prominently accessible
from the main navigation and not hidden.

% ─────────────────────────────────────────────────────────────────────────
\section{Model Governance}
% ─────────────────────────────────────────────────────────────────────────

\subsection{Sign-Off Status}

As of 2026-04-02, all nine sign-off areas are approved by the model owner.

\begin{center}
\begin{tabular}{llc}
\toprule
Area & Description & Status \\
\midrule
G1 & Data pipeline and panel integrity      & APPROVED \\
G2 & OLS fair-value specification           & APPROVED \\
G3 & SDE parameter estimation               & APPROVED \\
G4 & Monte Carlo simulation correctness     & APPROVED \\
G5 & Scoring engine calibration             & APPROVED \\
G6 & Validation and diagnostics pack        & APPROVED \\
G7 & UI disclosure and limitations          & APPROVED \\
G8 & Risk register completeness             & APPROVED \\
G9 & Reproducibility (replay deviation = 0) & APPROVED \\
\bottomrule
\end{tabular}
\end{center}

\subsection{Risk Register}

The risk register contains 26 items:
2 RESOLVED, 2 MONITORED, 5 PARTIALLY\_RESOLVED, 17 open/legacy.

Key resolved items as of 2026-04-02:

\begin{itemize}[leftmargin=1.8em]
  \item \textbf{R002 (RESOLVED):} Legacy scripts archived to
        \texttt{archive/legacy\_model/}.  The \texttt{src/model/} directory
        now contains 7 canonical files only.  No legacy imports remain in tests.
  \item \textbf{R004 (RESOLVED):} Score labels changed from Buy/Hold/Don't~Buy
        to Supportive/Mixed/Cautious; UI disclosure added;
        documented in \texttt{SCORING\_BACKTEST\_NOTE.md}.
\end{itemize}

\subsection{Model Owner Review Schedule}

Quarterly.  Next review: 2026-06-30.

\subsection{Deployment Restrictions}

\begin{govbox}[Deployment Restrictions]
This model \textbf{may not} be used for: (i)~individual property
valuations; (ii)~short-horizon trading decisions (horizon $<$24 months);
(iii)~regulatory capital calculations; or (iv)~any purpose requiring a
stress-testing model as defined under PRA~SS3/18.  These restrictions are
disclosed in the UI Limitations page, in this report, and in the risk
register.
\end{govbox}

% ─────────────────────────────────────────────────────────────────────────
\section{Limitations and Future Work}
% ─────────────────────────────────────────────────────────────────────────

\subsection{Current Limitations}

\begin{enumerate}[leftmargin=1.8em, label=\textbf{L\arabic*.}]
  \item \textbf{Regional averages only.}  The model operates on Land Registry
        regional average prices.  It is not applicable to individual properties,
        sub-regional markets, or property-type segments.
  \item \textbf{Interpolated annual series.}  Earnings and net housing additions
        are available at annual frequency and are linearly interpolated to monthly.
        Within-year timing precision is therefore limited.
  \item \textbf{Omitted-variable and regime risk.}  Despite HAC correction,
        omitted macro variables (credit conditions, planning policy shocks,
        stamp duty changes) and structural breaks remain genuine risks.
  \item \textbf{Scenario Lab is a perturbation tool.}  The interactive Scenario
        Lab applies coefficient-based overlays to existing Stage~5 outputs; it
        is not a full econometric re-run with new data.
  \item \textbf{Score compression.}  Composite scores reduce multi-dimensional
        distributional information to a single number.  Evidence-level
        classifications (Validated/Partial/Contextual) partially offset this.
  \item \textbf{CI coverage marginal.}  Pooled 1-step-ahead coverage of 78.7\%
        is below the 80\% pass threshold.  South West and Scotland require
        monitoring.
  \item \textbf{Income sign violation.}  The OLS income coefficient is negative
        (Frisch-Waugh-Lovell artefact); the model relies on the rent channel
        for income-affordability signal.
  \item \textbf{Earnings staleness.}  At the modelling date, earnings data
        are 13 months stale across all regions.
\end{enumerate}

\subsection{Future Extensions}

\begin{enumerate}[leftmargin=1.8em, label=\textbf{F\arabic*.}]
  \item Sub-regional and property-type granularity using ONS sub-regional HPI.
  \item Higher-frequency earnings proxies (e.g.\ PAYE real-time information)
        to reduce the 13-month staleness.
  \item Regime-switching extension replacing the fixed-$\kap$ OU process with
        a hidden Markov model for boom/bust regimes.
  \item Formal structural break test using the Andrews (1993) supremum-$F$
        procedure \cite{andrews1993}.
  \item Dynamic scenario weights incorporating macro-financial nowcasting signals
        rather than fixed judgement-based probabilities.
\end{enumerate}

% ─────────────────────────────────────────────────────────────────────────
\section{Conclusion}
% ─────────────────────────────────────────────────────────────────────────

This report has documented the UK Regional Housing Market Model~v0.2.0
from its probability-theoretic foundations through to production
validation and governance.  The principal achievements are:

\begin{enumerate}[leftmargin=1.8em]
  \item A complete, reproducible seven-stage pipeline from raw Land Registry
        and ONS data to an interactive Streamlit research dashboard.
  \item A semi-structural OLS fair-value layer ($n=452$, $R^2=0.9496$)
        with HAC standard errors and documented sign-violation disclosure.
  \item Region-level log-OU SDE parameters ($\kap \in [0.131, 0.825]$,
        $\sig \in [0.030, 0.120]$) estimated from AR(1) gap regressions
        and volatility blending.
  \item A log-OU simulation framework that outperforms GBM and zero-growth
        at the 12-month horizon: RMSE 0.04442 vs 0.04881 (GBM) and 0.05582
        (zero-growth), directional accuracy 80.6\% vs 79.0\% and 0.0\%.
  \item A bank-grade validation pack: replay deviation zero; subsample
        directional accuracy 64.7\% at 12 months (above 55\% threshold);
        CI coverage 78.7\% (MARGINAL, monitored); 100\% Ljung-Box and ARCH
        flagging mitigated by HAC standard errors.
  \item A full governance framework: 9/9 sign-off areas approved; 26-item
        risk register; score label governance; explicit deployment restrictions.
\end{enumerate}

The Euler-Maruyama log-OU process with a time-varying fair-value mean
path is the central modelling choice that makes the simulation
economically coherent and superior to GBM at planning-relevant horizons.
The model is suitable for systematic institutional research into UK
regional housing dynamics and long-horizon scenario analysis for REIT
portfolio management.

% ─────────────────────────────────────────────────────────────────────────
\begin{thebibliography}{9}

\bibitem{case1989}
  Case, K.~E.\ \& Shiller, R.~J.\ (1989).
  The efficiency of the market for single-family homes.
  \textit{American Economic Review}, 79(1), 125--137.

\bibitem{newey1987}
  Newey, W.~K.\ \& West, K.~D.\ (1987).
  A simple positive semi-definite, heteroskedasticity and
  autocorrelation consistent covariance matrix.
  \textit{Econometrica}, 55(3), 703--708.

\bibitem{kloeden1992}
  Kloeden, P.~E.\ \& Platen, E.\ (1992).
  \textit{Numerical Solution of Stochastic Differential Equations}.
  Springer, Berlin.

\bibitem{andrews1993}
  Andrews, D.~W.~K.\ (1993).
  Tests for parameter instability and structural change with
  unknown change point.
  \textit{Econometrica}, 61(4), 821--856.

\bibitem{vasicek1977}
  Vasi\v{c}ek, O.\ (1977).
  An equilibrium characterization of the term structure.
  \textit{Journal of Financial Economics}, 5(2), 177--188.

\end{thebibliography}

\end{document}
